\section{Hankel geometry forced by calibrated binary observations}\label{sec:binary-hankel}
We now return to the binary polynomial observation, rather than choosing
contrasts to realize a desired metric. Fix two distinct strictly positive
rank-one stochastic tables $M_a=u_av_a^{\mathsf T}$, $a=1,2$, with calibrated
channels. At a clock $T>D$ observe the four-cell law
\begin{equation}\label{eq:calibrated-binary}
 p(T)=\pi(T)M_1+(1-\pi(T))M_2,\qquad
 \frac{\pi(T)}{1-\pi(T)}=
 \frac{\alpha f_1(T)}{(1-\alpha)f_2(T)},\quad 0<\alpha<1.
\end{equation}
Both component polynomials have roots in $[0,D]$ and are positive at all
clocks. Fixed polynomial factors are permitted in either component.
The variable factors consist of $k\ge1$ separated interior double clusters
and $e\in\{0,1,2\}$ simple endpoints, with no shared variable cluster
between components. An endpoint is either $0$ or $D$, and occurs at most
once. Give a variable cluster the colour $\sigma_a=1$ or $-1$ according
as it belongs to component one or two. All other factors and the channels
are held fixed. Equal prescribed component degrees can be obtained by
fixed background factors; no coefficient is chosen from a target metric.

For an interior cluster at $r_i$ use its sum displacement $s_i$ and variance
$v_i$ through the factor $(T-r_i-s_i/2)^2-v_i$. The variance is nonnegative.
A lower endpoint has factor $T-z$, and an upper endpoint factor $T-D+z$,
where $z\ge0$. The log mixing odds is a free coordinate, as are the $k$
interior sums. Its use only specifies a smooth local coordinate in the
original rational model with mixing parameter $\alpha$.
The $n=k+e$ variances and inward endpoint shifts are retained.

Choose exactly $m=2k+e+1$ distinct clocks $D<T_1<\cdots<T_m$ and positive
exposures $\lambda_j$. Exposure means that the number of observations in
group $j$, divided by a common asymptotic scale, tends to $\lambda_j$.
The total exposure is not fixed to one in the cone theorem below. This is
an ordinary product of categorical experiments with its own Fisher metric.
At the centre put
\begin{equation}\label{eq:native-kappa}
 \kappa_j=\pi_j^2(1-\pi_j)^2
       \sum_{c=1}^4\frac{(M_{1c}-M_{2c})^2}{p_{jc}}>0.
\end{equation}
Order the retained sites as $\rho_1<\cdots<\rho_n$. Set
$\epsilon_a=-\sigma_a$ at an interior cluster or a lower endpoint and
$\epsilon_a=\sigma_a$ at an upper endpoint. Define
\begin{equation}\label{eq:binary-score}
 L_j=\left(1,\left(-\frac{\sigma_i}{T_j-r_i}\right)_{i=1}^k\right),
 \qquad
 D_{ja}=\frac{\epsilon_a}{(T_j-\rho_a)^{\mu_a}},
\end{equation}
where $\mu_a=2$ at an interior site and $1$ at an endpoint. The square
matrix $V=[L,D]$ has $m$ columns.

\begin{theorem}[Native finite moment cone]\label{thm:binary-hankel}
For the model \eqref{eq:calibrated-binary}, profiling the free sum and
mixing coordinates gives a positive retained information matrix $Q$.
Let $R=Q^{-1}$ and let $Z$ be the last $n$ rows of $V^{-1}$. Then
\begin{equation}\label{eq:native-inverse}
 R=\sum_{j=1}^m\frac{Z_{\cdot j}Z_{\cdot j}^{\mathsf T}}
                            {\lambda_j\kappa_j}.
\end{equation}
Its columns are fixed by the observation model. More explicitly, set
\[
 h(T)=\prod_{i=1}^k(T-r_i)^2\prod_{b\text{ endpoint}}(T-b),\qquad
 P(T)=\prod_{j=1}^m(T-T_j),\quad p_*(T)=\prod_{a=1}^n(T-\rho_a),
\]
\[
 h_a(\rho_a)=\left.\frac{h(T)}{(T-\rho_a)^{\mu_a}}\right|_{T=\rho_a},
 \quad d_a=-\frac{P(\rho_a)}{\epsilon_a h_a(\rho_a)},\quad
 \xi_j=\frac{h(T_j)}{P'(T_j)}.
\]
Then
\begin{equation}\label{eq:cauchy-columns}
 Z_{aj}=\frac{\xi_j d_a}{T_j-\rho_a}.
\end{equation}
Let $E$ be the coefficient matrix whose $a$th row is the polynomial
$p_*(T)/(T-\rho_a)$ in the basis $1,T,\ldots,T^{n-1}$, and put
$J_* =\diag(d_a)E$. It is invertible. With
$w(T)=(1,T,\ldots,T^{n-1})^{\mathsf T}$, the entire native image is
\begin{equation}\label{eq:hankel-cone}
 \begin{split}
 R&=J_*H(\omega)J_*^{\mathsf T},\\
 H(\omega)&=\sum_{j=1}^m\omega_j w(T_j)w(T_j)^{\mathsf T},\qquad
 \omega_j=\frac{\xi_j^2}{p_*(T_j)^2\lambda_j\kappa_j}>0.
 \end{split}
\end{equation}
As the positive exposures vary, $H$ ranges over the relative interior of
the finite moment cone generated by these clock tensors. This cone has
linear span equal to the full Hankel subspace and dimension $2n-1$.
Consequently the native $Q$-image is a smooth $(2n-1)$-dimensional image
under matrix inversion. These assertions hold for every separated clock
and root configuration in the stated class, not merely generically.
\end{theorem}
\begin{proof}
Differentiating the log odds of \eqref{eq:calibrated-binary} gives exactly
\eqref{eq:binary-score}. The probability derivative in a log-odds direction
is $\pi_j(1-\pi_j)(M_1-M_2)$, so its scalar information is
\eqref{eq:native-kappa}. The full information is
$I=V^{\mathsf T}\diag(\lambda_j\kappa_j)V$.
To prove $V$ invertible, multiply a linear combination of its rational
columns by $h(T)$. The result is a polynomial of degree at most $m-1$.
If it vanishes at all $m$ clocks it is identically zero. Its constant
part at infinity and its simple and double principal parts at each site
then force every coefficient to vanish. Thus $I>0$. The lower-right
block of $I^{-1}$ is the inverse of the nuisance Schur complement, giving
\eqref{eq:native-inverse}.

For explicit recovery, prescribe arbitrary clock log-odds increments $y_j$.
The polynomial $g(T)=h(T)(V\theta)(T)$ is reconstructed by Lagrange
interpolation as
\[
 g(T)=\sum_j h(T_j)y_j\frac{P(T)}{(T-T_j)P'(T_j)}.
\]
At a retained site, its highest principal-part coefficient is
$g(\rho_a)/(\epsilon_a h_a(\rho_a))$. This proves
\eqref{eq:cauchy-columns}. The polynomials $p_*(T)/(T-\rho_a)$ are a
basis of the polynomials of degree at most $n-1$, as evaluation at the
$n$ sites immediately shows. Hence $J_*$ is invertible and
\[
 Z_{\cdot j}=\frac{\xi_j}{p_*(T_j)}J_*w(T_j).
\]
Substitution proves \eqref{eq:hankel-cone}. Positive exposures give all
positive $\omega$, independently. A linear map sends the positive orthant
onto the relative interior of its finitely generated cone: one direction
follows by openness onto the span, and for the converse subtract a small
positive multiple of the sum of all generators from a relative-interior
point before representing the remainder with nonnegative coefficients.

The entries of $w(T)w(T)^{\mathsf T}$ have powers $T^s$ for
$0\le s\le2n-2$. Since $m=2k+e+1\ge2n-1$ exactly when $e\le2$, the
Vandermonde matrix at the distinct clocks has rank $2n-1$. The generators
therefore span precisely the Hankel space. Their positive combinations
are positive definite, because the vectors $w(T_j)$ span $\R^n$.
Inversion is a diffeomorphism on the positive cone, which proves the
last dimension assertion.
\end{proof}

\begin{corollary}[Secants, dual queries, and fixed budgets]\label{cor:hankel-secants}
Let $\mathcal V_*=J_*\{\text{Hankel matrices}\}J_*^{\mathsf T}$ and
$s=2n-1$. The normalized nonzero secants of the native inverse-information
image are exactly the Frobenius unit sphere of $\mathcal V_*$. Recovering
an unknown $R$ from exact dual values $u^{\mathsf T}Ru$, with $u$ in any
specified cone with nonempty interior, requires and admits exactly $s$
queries, both for fixed queries and for deterministic adaptive queries.
The exposures are not supplied as additional numerical input to this
oracle problem.

For the information matrices themselves use normalized coordinates
$\widetilde Q=J_*^{\mathsf T}QJ_*=H^{-1}$. A symmetric unit matrix $D$
is a normalized secant of this image if and only if there exist
$H_0,H_1$ in the relative-interior moment cone and $a>0$ with
\begin{equation}\label{eq:inverse-secants}
 H_1DH_0=a(H_0-H_1).
\end{equation}
This is a finite polynomial presentation in positive clock weights,
$a$, and the entries of $D$. Physical-coordinate secants are obtained by
$D\mapsto J_*^{-\mathsf T}DJ_*^{-1}$ and renormalization.
The sharp dual-query count is not asserted for the different oracle
$v^{\mathsf T}Qv$.

If the total exposure is fixed to one, the exact image is obtained by
adding, to \eqref{eq:hankel-cone}, the condition
\begin{equation}\label{eq:fixed-budget}
 \sum_j\frac{\xi_j^2}{p_*(T_j)^2\kappa_j\omega_j}=1.
\end{equation}
No unconstrained-cone dimension or secant assertion is imposed on that
budget section.
\end{corollary}
\begin{proof}
The image of $R$ is relatively open in $\mathcal V_*$. A relative ball
around any image point realizes every nonzero direction as a secant;
all differences lie in that span. Quadratic evaluations on a cone with
interior separate $\mathcal V_*$, and $s$ of them form a basis of its
dual. For fewer than $s$ adaptive evaluations at an interior point,
a nonzero matrix in their common kernel gives two sufficiently small
opposite perturbations in the image with identical transcripts. This
proves the exact count without continuity assumptions on the algorithm.
Multiplying the identity
$D=a(H_1^{-1}-H_0^{-1})$ by $H_1$ and $H_0$ proves both directions of
\eqref{eq:inverse-secants}. The formula for $\omega_j$ gives
\eqref{eq:fixed-budget} directly.
\end{proof}
For $n=2$, the dimension is three, consistent with the retained
 two-variance phenomenon. For $n\ge3$, the exact dual image has Hankel
relations rather than arbitrary symmetric entries. The coordinate
congruence and the distinction between a tensor and its inverse are
part of the classification, not changes in the observation metric.

\begin{theorem}[Native endpoint visibility]\label{thm:native-visibility}
In the subclass in which all variable factors belong to one component,
partition $Q$ into interior variances and inward endpoints. For $e=1$
and $k\ge2$ the representation is fully visible. For $e=2$ and $k\ge3$
it is fully visible. Thus its minimal endpoint dimension and its full
minimal fibre are reconstructed by Theorem~\ref{thm:reconstruct}, for
every positive exposure vector.

For $e=2,k=2$ it has exactly three cells, with the empty cell missing;
its entire minimal fibre is \eqref{eq:fork-fibre}. For $k=1$ and
$e=1$ or $2$, all endpoints are active on the positive retained ray;
the profile is a single positive quadratic form, with minimal endpoint
dimension zero. These conclusions are intrinsic to this binary family.
\end{theorem}
\begin{proof}
When the variable factors have one colour, all $d_a$ have the same sign.
Indeed, the squared interior factors in $h_a$ are positive, and each
upper-endpoint sign reversal compensates the negative factor contributed
by that endpoint at an interior site. Directly checking the possible
sets $\varnothing,\{0\},\{D\},\{0,D\}$ gives a common sign.
The matrix $R$ is consequently a positive diagonal congruence of a
Cauchy Gram matrix
\[
 \left(\frac1{T_j-\rho_a}\right)_{a,j}
 \diag(\gamma_j)
 \left(\frac1{T_j-\rho_a}\right)_{a,j}^{\mathsf T},\qquad\gamma_j>0.
\]
Every minor of size $h$ of the rectangular Cauchy matrix has the same
nonzero sign $(-1)^{h(h-1)/2}$, by the Cauchy determinant formula.
Cauchy--Binet therefore makes every minor of $R$ strictly positive.
Jacobi's complementary-minor identity shows that
$Q'=\Sigma Q\Sigma$, $\Sigma_{aa}=(-1)^a$, is also strictly totally
positive. We use only these displayed determinant identities, not a
separate hypothesis of total positivity.

For one endpoint the entries in the cross column alternate signs along
the ordered interior sites. With at least two sites they have both
signs, so a strictly positive vector annihilates the transpose of the
cross block. Proposition~\ref{prop:positive-span} proves visibility.
For two endpoints at the extreme sites, after fixed column sign changes
the interior rows of $B$ have alternating signs times positive vectors
$(a_i,b_i)$. The ratios $a_i/b_i$ are strictly monotone, because every
ordered $2\times2$ minor of $Q'$ is positive. For three consecutive
rows, the middle positive vector lies strictly in the cone of the two
outer ones, and its sign is opposite. Hence these rows admit a strictly
positive linear dependence. Additional rows can receive sufficiently
small positive weights; a bounded right inverse of the three-row block
corrects their residual while keeping the first three weights positive.
Thus $B$ has rank two and $B^{\mathsf T}h=0$ for some $h>0$. The same
proposition proves full visibility for all $k\ge3$.

For $e=2,k=2$ order the sites as endpoint, interior, interior, endpoint.
Then
\[
 B=\begin{pmatrix}-a&b\\c&-d\end{pmatrix},\qquad
 a,b,c,d>0,\qquad ad>bc.
\]
The last inequality is a minor of $Q'$. The empty-cell conditions would
require $x_2/x_1\ge a/c$ and $x_2/x_1\le b/d$, which are incompatible.
The endpoint block has negative off-diagonal, so its inverse is entrywise
positive. Taking $b/d<x_2/x_1<a/c$ gives $B^{\mathsf T}x<0$ and a
strict full-active point. At $x=(1,0)$ the first-only inactive multiplier
is
$(Q_{11}Q_{24}-Q_{14}Q_{21})/Q_{11}>0$; at $x=(0,1)$ the second-only
multiplier is $(Q_{31}Q_{44}-Q_{14}Q_{34})/Q_{44}>0$.
These are positive signed minors, and the strict inequalities persist
on nearby interior rays. Exactly three cells occur, so
Theorem~\ref{thm:three-cells} applies.

If $k=1$, block inversion gives
$-C^{-1}B^{\mathsf T}=R_{E,I}R_{I,I}^{-1}$, which is entrywise positive.
Every endpoint is therefore active on $x>0$. Completing the square gives
a single positive Schur-complement quadratic, proving the last assertion.
\end{proof}

For several component colours, Theorem~\ref{thm:binary-hankel} and
Corollary~\ref{cor:hankel-secants} remain unchanged. After conjugation by
$\diag(\operatorname{sign}d_a)$, $R$ is strictly totally positive; this
specifies the signed native class. In particular the physical signs are
not discarded when testing endpoint visibility. The all-one-colour
subclass above is a naturally calibrated submodel of the original binary
observations, not a realization obtained by first prescribing $Q$.
